% \subsection{Train-Test Split}
% \noindent - Standard split: split 0 is the held-out test set, splits 1–4 form the training set\\
% - The contamination problem: explain the protein cluster (pc) column -> a protein's cluster may include proteins from phages in different folds. The models that produced the scores were trained on protein clusters, so a protein from a training-fold phage may share a cluster with a test-fold protein, introducing subtle leakage\\
% - The masking fix: the file all\_random\_pc\_splits.pkl encodes for each protein and each (host, function) model which fold that protein was assigned to. Setting scores to NaN where the protein was not in the test set of the relevant model, before aggregation, removes the contaminated predictions. The existing NaN-aware mean pooling then handles these naturally -> this is the "missing value" framing\\


\section{Task Specific Filtering and Data Partitioning}
\label{sec4.3: contamination fix}

The main requirement for the finite sample conformal validity (Theorem~\ref{thm: marginal coverage}) is the 
exchangeability of the calibration samples $S_2$ and test samples $S_{\text{test}}$. In our dataset, the cross validation 
partition is defined at the phage level ($f_p \in \{0,1,2,3,4\}$), whereas the underlying base scoring models 
were trained at the protein-cluster (\texttt{pc}) level. A single protein cluster can span phages assigned to different genome 
folds, therefore a base model scoring protein $p$ for category $k$ may have observed homologous proteins from $p$'s cluster 
during training.

Formally, let $F_{p,k} \in \{0,1,2,3,4\}$ denote the fold held out when training the base model for feature $k$. A naive 
aggregation which uses scores where $F_{p,k} \neq f_p$ introduces optimistic bias, by carrying training information into test 
evaluations. This breaks the exchangeability assumption that is a requirement of Theorem~\ref{thm: marginal coverage}.

The underlying prediction models for host level and gram type classification differ in structure, hence we correct for data 
contamination at the appropriate granularity for each task:

For \textbf{host level classification}, each of the $M = 2\,160$ host function models were trained on an independent 
five fold split of the protein cluster space. Let $F_{p,k} \in \{0, 1, 2, 3, 4\}$ denote the held out fold identifier 
for protein $p$ under the model for feature $k$. To prevent homology leakage across splits, we enforce an out of fold 
masking operator prior to phage aggregation
\begin{equation}
\label{eq: contamination-mask}
s_{p,k} \leftarrow \text{NA} \quad \text{if } F_{p,k} \neq f_p.
\end{equation}

We replace the in fold predictions with $\texttt{NA}$, then the NaN-aware mean pooling operator \eqref{eq: nan-mean-pooling} 
naturally excludes the contaminated scores without introducing imputation bias. This allows us to utilize the full  dataset across 
all five folds without discarding data:
\begin{equation*}
D_{\text{test}} = \{i : \text{split}_i = 0\}, \qquad D_{\text{val}} = \{i : \text{split}_i = 1\}, \qquad D_{\text{train}} = 
\{i : \text{split}_i \in \{2,3,4\}\}.
\end{equation*}

Here, $D_{\text{train}}$ provides us with the the shape-discovery and size-scaling calibration sets ($S_1, S_2$), 
$D_{\text{val}}$ is used to tune envelope hyperparameters (e.g., bin count $NB$, directional resolution $M$, and kernel smoothing 
$\kappa$), and $D_{\text{test}}$ is held out for final evaluation.

\newpage
For \textbf{gram type classification}, the $K = 40$ functional category models used for gram type classification do not 
have per column split assignment tables. As a conservative, unambiguous fallback, we restrict the gram type dataset strictly to 
phages in Fold 0 ($\text{split}_i = 0$). Because Fold 0 phages are consistently held out at the genome cluster level, all 
associated protein scores are unbiased by construction. Within this Fold 0 subset, we partition the phages $40/30/30$ into shape 
discovery ($S_1$), size-scaling ($S_2$), and test ($S_{\text{test}}$) sets following the standard CSA calibration split.

In both cases, $M_{p,k}$, which is the indicator determining which protein level scores are treated as valid input 
to the pooling operator of Equation~\eqref{eq: nan-mean-pooling}, reduces to whichever correction each task's models 
support. This is the exact fold match test of Equation~\eqref{eq: contamination-mask} for host level classification, or 
membership in the Fold-$0$ subset for gram type classification.


% \subsubsection{Valid Prediction Filter}
% \label{sec4.3.2: prediction filter}

% \noindent To ensure that only clean, out-of-fold scores are used during phage level aggregation, we define the binary mask $M_{p,k} 
% \in \{0,1\}$ as a product of four simple checks:
% \begin{equation}
% \label{eq: validity-mask-def}
% M_{p,k} = \mathbf{1}(F_{p,k} = f_p) \cdot C_{\text{clean}}(p) \cdot Q(p,k) \cdot D(p,k),
% \end{equation}
% \noindent where:
% \begin{enumerate}[label=(\roman*)]
%     \item $\mathbf{1}(F_{p,k} = f_p)$ is the is the cross-validation fold-match indicator defined in \eqref{eq: contamination-mask}. 
%           It ensures the prediction comes from an out-of-fold model \eqref{eq: contamination-mask}.
%     \item $C_{\text{clean}}(p) \in \{0,1\}$ removes host genomic contamination (e.g., $>95\%$ identity to non-viral host regions).
%     \item $Q(p,k) \in \{0,1\}$ checks sequence quality (e.g., length $L_p \ge 30\text{ aa}$ and $E\text{-value} < 10^{-3}$).
%     \item $D(p,k) \in \{0,1\}$ verifies domain applicability, marking $s_{p,k} = \text{NA}$ if $p$ falls outside the functional 
%           domain of category $k$ .
% \end{enumerate}
% \noindent If any check fails ($M_{p,k} = 0$), the score $s_{p,k}$ is set to $\texttt{NA}$ and ignored during pooling 
% \eqref{eq: nan-mean-pooling}. This ensures that every score reaching the calibration step is unbiased and exchangeable.







